\section{2D Segmentation and 3D Reconstruction}
\markright{\thesection ~~~ Methodology1}
\label{method1}

The first approach to the problem was a computationally lightweight methodology, inspired by the work of Rizvandi et al.\ \cite{rizvandi} on the detection of overlapped nematode worms. Despite the apparent difference in application domain, the two problems share a fundamental structural challenge: given a skeleton extracted from a binary image, deciding which segments that meet at a junction point belong to the same continuous object and which ones do not. In both cases, the core difficulty lies in resolving ambiguous intersections between elongated, curvilinear structures, a problem that is geometrically equivalent whether the objects in question are biological organisms or glass fibers embedded in a polymer matrix. 

\textcolor{red}{Adicionar visão geral do processo de segmentação e reconstrução.}

To detect and separate individual \textit{C.\ elegans} nematode worms in population images containing overlapping specimens, Rizvandi et al.\ \cite{rizvandi} proposed a four-stage pipeline based on skeleton analysis. In the pre-processing stage, the input image is smoothed with a Gaussian kernel and binarized using Otsu's global thresholding method, followed by morphological erosion and hole-filling operations to produce a clean binary mask. A morphological skeletonization is then applied, and spurious short branches resulting from residual noise are pruned by splitting the skeleton at junction pixels (defined as pixels with more than two neighbors) and discarding branches below a minimum length threshold. In the skeleton analysis stage, each junction pixel is removed, decomposing the skeleton into a set of isolated branches. For each connection pixel at the boundary of a branch, a local orientation angle is computed by averaging the successive angular differences across a window of neighboring pixels along that branch. In the angle analysis and tree construction stage, these local angles are used to decide which branches should be joined: starting from each free endpoint, a tree is grown by iteratively appending the neighboring branch whose orientation is most consistent with the current trajectory, accepting the merge only if the angular difference between the two branches falls below a user-defined threshold. This greedy angle-based merging reconstructs the continuous path of each individual worm through the junction points without requiring any prior knowledge of the objects' positions. Finally, a post-processing stage filters out spurious detections by retaining only those reconstructed paths whose length falls within a statistically determined interval, derived from a reference set of isolated individuals.
% Applied to a dataset of 147 isolated and overlapping worms, the method achieved an accuracy of 81.43\%, outperforming the semi-automatic reference method it was compared against, which required manual marking of head and tail positions and reached only 55.33\% accuracy on the same dataset.

\subsection{Slicing Axis Selection and Pre-processing}

Analysing the 3D micro-CT image (Figure~\ref{fig:raw}), the first decision to be made was the choice of slicing axis --- that is, the direction along which the volumetric image would be iteratively sectioned into 2D layers, each of which would be individually segmented before being stacked back to reconstruct a fully segmented 3D representation. A previous attempt by a researcher at the Computational Intelligence Laboratory (LITC) applied the OpenFiberSeg method \cite{sosa-rey} by slicing along the X axis, producing layers in the \textit{ZY} plane in which most fibers appear running parallel to the cross-section (left image in Figure~\ref{fig:sections}). The results obtained with that approach were unsatisfactory, for the reasons discussed in the previous chapter. The present work therefore adopts a different perspective, slicing the volume along the Z axis and producing \textit{XY} layers (right image in Figure~\ref{fig:sections}). This choice is justified by the expectation that most fibers are preferentially aligned with the X axis; consequently, each individual fiber should appear in only a small number of consecutive \textit{XY} layers --- ideally on the order of its mean diameter expressed in pixels. With each fiber spanning fewer layers, the subsequent 3D reconstruction becomes considerably more tractable. The Z axis was preferred over the Y axis as the slicing direction because it yields a smaller total number of layers, thereby reducing the computational cost of the reconstruction step.

\begin{figure}
    \centering
    \includegraphics[width=0.9\linewidth]{Metodologia/Figuras/eixos.png}
    \caption{Micro-CT image with reference axes.}
    \label{fig:raw}
\end{figure}

\begin{figure}
    \centering
    \includegraphics[width=0.7\linewidth]{Metodologia/Figuras/sections.png}
    \caption{Representative layer examples for each slicing axis.}
    \label{fig:sections}
\end{figure}

In the pre-processing stage, the 3D image was binarized using Otsu's global thresholding method \cite{otsu}, which automatically determines an intensity threshold that maximally separates foreground (fiber) and background voxels. The resulting binary image was then subjected to a morphological opening filter (a sequential erosion followed by a dilation) which smooths the fiber boundaries and suppresses spurious small protrusions introduced by the binarization process. Alternative filters, including median filtering and morphological closing (hole filling), were evaluated but discarded, as they caused artificial thickening of the fiber cross-sections and consequently increased the degree of apparent overlap between adjacent fibers.

With the filtered binary image in hand, a morphological skeletonization was applied \textcolor{red}{(REF IMAGEM DO ESQUELETO)}, reducing each fiber cross-section to a single-pixel-wide medial axis. Skeletonization is a critical pre-processing step in this pipeline: if the full binary image were used directly, regions where two or more fibers overlap would appear as extended multi-pixel blobs, making the subsequent disambiguation between touching fibers considerably more difficult. By contrast, the unifilar skeleton collapses each such overlapping region into a single pixel or a very small cluster of pixels, greatly simplifying the structural analysis of fiber junctions.

\subsection{2D Angle Analysis and Pruning}

The pruning stage constitutes the core of the 2D segmentation step, in which the algorithm decides, at each junction point of the skeleton, which pairs of converging branches belong to the same continuous fiber and which ones originate from distinct fibers crossing one another. To formalize this decision process, it is necessary to first define the structural elements of the skeleton.

A \textbf{junction pixel} is any skeleton pixel that has three or more immediate neighbors also belonging to the skeleton, indicating a point where multiple fiber paths converge. An \textbf{endpoint pixel} is a skeleton pixel with exactly one neighbor, marking the terminal extremity of a fiber or branch. A \textbf{branch} is then defined as a connected chain of skeleton pixels linking two such special pixels. Three branch types are distinguished: a \textit{JP--JP branch} connects two junction pixels; a \textit{JP--EP branch} connects a junction pixel to an endpoint; and an \textit{EP--EP branch} connects two endpoints, corresponding to an isolated fiber segment that does not intersect any other structure.

Once all branches have been mapped, a filtering step is applied in which each branch is assessed according to its type and its length, measured as the number of pixels composing the branch. Branches of type JP--EP or EP--EP whose length falls below a minimum threshold of 5 pixels \textcolor{red}{(checar código e explicar escolha, acho que faço o filtro após a poda, o que pode ser ruim)} are removed by setting their pixels to zero, as they are considered to be artefacts of the skeletonization or residual noise rather than meaningful fiber segments. JP--JP branches, however, are not filtered by length, since removing short inter-junction segments could erroneously sever a genuine fiber at a location where two junction points happen to lie in close proximity along its axis.

The outcome of this filtering step is a cleaner set of 2D skeleton images with fewer spurious fragments, which facilitates the subsequent decision-making at junction points \textcolor{red}{(verificar se o OBJ esqueleto é recalculado depois da filtragem, se é que a filtragem é feita pré)}. The segmentation logic then iterates over all junction pixels in each layer. At each junction pixel, the local orientation angle of every incident branch with respect to the Y axis is computed. To ensure that this angle accurately reflects the fiber's trajectory in the immediate vicinity of the junction, rather than its overall direction across the full branch length, the angle is calculated using a fixed local window: for branches longer than 12 pixels, the angle is determined by the line connecting the junction pixel to the 12th pixel along the branch; for shorter branches, the endpoint is used instead. This local windowing strategy was adopted after observing that computing the angle from the junction point to the branch endpoint produced unreliable results for curved fibers, as the global direction of a long curved branch can diverge significantly from its actual local orientation near the junction, leading to erroneous collinearity decisions, and it's threshold of 12 pixels was chosen empirically. These locally computed angles are then compared pairwise across all branches meeting at the junction: the pair whose angles sum to a value closest to 180° are considered collinear and assigned to the same fiber, provided that their sum is no less than 45°. This lower bound was chosen empirically and acts as a sharpness filter, a junction angle below 45° would imply an abrupt change in fiber direction that is physically incompatible with the behavior of a rigid glass fiber, and is therefore interpreted as evidence of two distinct fibers crossing rather than a single fiber continuing through the junction. The remaining branches at that junction are interpreted as belonging to distinct fibers and are pruned by setting their first pixel beyond the junction point to zero, effectively disconnecting them from the junction and separating the crossing fiber paths \textcolor{red}{(se forem 4 corto dois que podem estar colineares também, sinal de "+")}.

This decision criterion relies on the assumption that fibers follow locally linear trajectories in the vicinity of each junction point. This assumption is physically grounded in the nature of the material system and its manufacturing process. Since the composite samples analyzed in this work were produced by thermoplastic injection molding \cite{aguiar-base}, the fibers are expected to behave as rigid rods throughout processing, as described by the Folgar--Tucker model \cite{folgar1984}. Under the high shear rates characteristic of injection molding, short glass fibers align preferentially along the flow direction and exhibit locally linear trajectories, a behavior that has been experimentally confirmed for this material system through micro-CT fiber orientation measurements \cite{aguiar-base}.

\subsection{Fiber Labeling}

After pruning the skeleton at each junction point, the resulting disconnected branches can be individually labeled, with each label corresponding to a distinct fiber skeleton segment. Ideally, each physical fiber will produce exactly one branch per layer, which is consistent with the nature of glass fibers: since they do not bifurcate or split during the injection molding process, no branching topology is expected in their skeletal representation.

The individual branches are identified and labeled using connected component analysis \textcolor{red}{(colocar REF e verificar código)}, a standard image processing operation that assigns a unique integer label to each set of pixels that are spatially connected to one another and disconnected from all others. Following the pruning step, each connected component in the skeleton corresponds to a single branch, and its pixels are assigned a common label value that serves as its identifier throughout the subsequent processing stages.

With all skeleton branches properly labeled, it is necessary to propagate those labels to the full binary fiber mask, so that every foreground pixel in the binary image is associated with a specific fiber identity rather than merely with the fiber class as a whole. This propagation is performed using the watershed transform \textcolor{red}{(colocar REF e verificar código)}, a morphological algorithm that treats the grayscale image as a topographic surface and floods it from a set of seed markers. In this context, the labeled skeleton pixels serve as the seed markers, and the watershed algorithm assigns each binary foreground pixel the label of its nearest skeleton branch, effectively partitioning the binary fiber regions according to their proximity to each skeletal segment. The result is a fully labeled 2D binary image in which every fiber pixel carries the identifier of the fiber branch it belongs to.

\subsection{3D Reconstruction}

Although the pruning and labeling steps described above operate on individual 2D layers, each physical fiber is a three-dimensional object that spans multiple consecutive cross-sections. Since each layer is processed independently, there is no guarantee that the same fiber receives the same label across adjacent layers. The 3D reconstruction step addresses precisely this problem: it establishes correspondences between labeled fiber segments across consecutive layers, propagating and unifying labels so that all cross-sections belonging to the same physical fiber share a common identifier.

Before adopting this layer-to-layer matching strategy, an alternative approach was considered: computing the length and orientation of each 2D labeled branch directly from the cross-sectional plane and using these measurements as proxies for the true 3D fiber geometry. This approach was discarded for a fundamental reason. A fiber appears in multiple consecutive layers not only because of its three-dimensional extent in space, but also because of its angular relationship with the slicing plane. A fiber that runs obliquely with respect to the \textit{XY} slicing plane will intersect a larger number of layers than one running perpendicular to it, and its cross-sectional projection will appear as an elongated elliptical region rather than a circular one. Consequently, the orientation and length derived from any single 2D cross-section would not be representative of the true 3D fiber geometry, and aggregating these measurements across layers without first establishing fiber identity would introduce systematic errors in both length and orientation estimates. The 3D reconstruction step is therefore a necessary precondition for any meaningful metric computation.

\textcolor{red}{Acabar de explicar a reconstrução 3D.}